\documentclass[11pt,a4paper]{article}

\usepackage[T1]{fontenc}
\usepackage[utf8]{inputenc}
\usepackage{lmodern}
\usepackage{microtype}
\usepackage[a4paper,margin=29mm,headheight=14pt]{geometry}
\usepackage{amsmath,amssymb,amsthm,mathtools}
\usepackage{xcolor}
\usepackage{fancyhdr}
\usepackage[colorlinks=true,linkcolor=blue!45!black,
  citecolor=blue!45!black,urlcolor=blue!50!black]{hyperref}

\definecolor{paperblue}{RGB}{28,66,102}
\hypersetup{
  pdftitle={Forgetting a Simple Root of a Binary Cubic},
  pdfauthor={Independent exact audit},
  pdfsubject={A coordinate-free model for a noninjective Keller map in dimension three},
  pdfkeywords={Jacobian conjecture, Keller map, binary cubic, resultant, etale morphism}
}

\pagestyle{fancy}
\fancyhf{}
\fancyhead[L]{\small\color{paperblue}Marked roots of binary cubics}
\fancyhead[R]{\small\color{paperblue}Coordinate-free construction}
\fancyfoot[C]{\thepage}
\renewcommand{\headrulewidth}{0.35pt}

\newtheorem{theorem}{Theorem}[section]
\newtheorem{proposition}[theorem]{Proposition}
\newtheorem{lemma}[theorem]{Lemma}
\newtheorem{corollary}[theorem]{Corollary}
\theoremstyle{remark}
\newtheorem{remark}[theorem]{Remark}

\DeclareMathOperator{\Sym}{Sym}
\DeclareMathOperator{\Res}{Res}
\DeclareMathOperator{\Spec}{Spec}
\newcommand{\A}{\mathbb A}
\newcommand{\C}{\mathbb C}
\newcommand{\Gm}{\mathbb G_{\mathrm m}}
\newcommand{\PP}{\mathbb P}
\newcommand{\dd}{\mathrm d}

\title{\color{paperblue}\bfseries
  Forgetting a Simple Root of a Binary Cubic\\[2mm]
  \large A coordinate-free model for a noninjective Keller map in dimension three}
\author{Independent exact audit}
\date{20 July 2026}

\begin{document}

\maketitle

\begin{abstract}
We describe intrinsically a polynomial map
$F\colon\A^3_{\C}\to\A^3_{\C}$ with constant Jacobian determinant
$-2$ and geometric degree three.  Let $W$ be a two-dimensional complex
vector space, let $H$ be a tangential affine slice of the space of binary
cubics, and mark one simple projective root of a cubic in $H$.  The resulting
simple-root incidence variety is, exceptionally, isomorphic to $\A^3$.
Under this identification, $F$ is the morphism which forgets the marked
root.  Resultant normalization gives a coordinate-free proof of etaleness;
the three roots of a general cubic explain the generic degree.  We prove the
affine-space identification, recover the explicit polynomial formula, and
interpret the discriminant as the precise locus of nonproperness.
\end{abstract}

\paragraph{Provenance.}
The explicit coordinate map studied here was publicly announced by Levent
Alp\"oge in an X post of 20 July 2026 \cite{Alpoge}; that post credits Akhil
for asking about the problem and Fable for working during the World Cup final.
Later that day Andy Jiang (@davikrehalt) posted the projective symmetric-product
formulation \cite{Jiang}, explicitly presenting it in the post as GPT output.
Aaron Lou subsequently gave a factorization--resultant derivation and a global
affine chart \cite{Lou}.  The purpose of this note is to make the equivalence
of these presentations precise and to supply a detailed intrinsic proof.

\section{The intrinsic construction}

Throughout, varieties are over $\C$.  Let $W$ be a two-dimensional vector
space and fix a nonzero alternating form
$\omega\in\bigwedge^2W^\vee$.  Choose a tensor
$\tau\in\Sym^3W$ of multiplicity type $(2,1)$: after extension to a basis it
has the form $u^2v$ with $u,v$ independent.  Equivalently, $[\tau]$ lies on
the tangent developable of the Veronese cubic

\[
  \nu_3\colon\PP(W)\longrightarrow\PP(\Sym^3W)
\]

but not on the Veronese curve itself.  Rescale $\tau$ once and for all, and
put

\[
  H_\tau=
  \left\{C\in\Sym^3W^\vee:\langle\tau,C\rangle=1\right\}.
\]

This is an affine hyperplane in a four-dimensional vector space, hence
$H_\tau\simeq\A^3$.  Consider the open simple-root incidence

\[
  X_\tau=
  \left\{(C,[r])\in H_\tau\times\PP(W):
  C(r)=0\ \text{and}\ \dd C_r\ne0\right\}.
\]

Our main statement is the following.

\begin{theorem}[Marked-root model]\label{thm:main}
The variety $X_\tau$ is isomorphic to $\A^3$.  Under this identification,
the forgetful morphism
\[
  \pi\colon X_\tau\longrightarrow H_\tau,
  \qquad (C,[r])\longmapsto C,
\]
is an etale polynomial map of affine three-spaces, generically of degree
three.  Its fiber over $C$ is naturally the reduced set of simple projective
roots of $C$.
\end{theorem}

\begin{remark}\label{rem:choices}
The construction is coordinate-free but not choice-free.  Its intrinsic data
are $(W,\omega,\tau)$ together with the chosen nonzero affine levels.  The
identifications of $X_\tau$ and $H_\tau$ with $\A^3$ are not canonical.
Accordingly, the statement that $\pi$ is etale is intrinsic, whereas the
numerical value $-2$ of its Jacobian depends on the chosen affine volume
coordinates.
\end{remark}

\subsection{The projective symmetric-product model}

The construction above is the affine normalization of the projective picture
posted by Jiang \cite{Jiang}.  Identify effective degree-three divisors on
$\PP(W)$ with projective binary cubics,

\[
  \Sym^3\!\bigl(\PP(W)\bigr)
  \simeq \PP\!\bigl(\Sym^3W^\vee\bigr),
\]

and let

\[
  \mu\colon \PP(W)\times\Sym^2\!\bigl(\PP(W)\bigr)
  \longrightarrow \Sym^3\!\bigl(\PP(W)\bigr),
  \qquad (p,D)\longmapsto p+D
\]

forget which root was marked.  Write $R$ for its ramification divisor and

\[
  H_\infty=\left\{[C]:\langle\tau,C\rangle=0\right\}.
\]

The small diagonal $\{3p:p\in\PP(W)\}$ is a twisted cubic.  Because $\tau$
has type $(2,1)$, the restriction of the defining functional of $H_\infty$
to this twisted cubic has zero divisor $2p+q$ with $p\ne q$.  Thus
$H_\infty$ is tangent but not osculating to the small diagonal.  Conversely,
every such hyperplane arises this way and all of them are equivalent under
$\operatorname{PGL}(W)$.

There are natural isomorphisms

\[
  \Sym^3\!\bigl(\PP(W)\bigr)\setminus H_\infty
  \xrightarrow{\sim}H_\tau,
  \qquad [C]\longmapsto
  \frac{C}{\langle\tau,C\rangle},
\]

and

\[
  \left(\PP(W)\times\Sym^2\!\bigl(\PP(W)\bigr)\right)
  \setminus\left(R\cup\mu^{-1}(H_\infty)\right)
  \xrightarrow{\sim}X_\tau.
\]

Under these isomorphisms, $\mu$ becomes $\pi$.  Indeed, deleting $R$ says
exactly that the marked root is simple; the two unmarked roots are still
allowed to coincide.  Lemma~\ref{lem:gauge} below gives a further affine lift:
writing the marked cubic as $C=\ell q$ and imposing
$\Res_\omega(\ell,q)=1$ fixes the remaining factor-rescaling gauge.  Hence
Jiang's projective restriction, the normalized factor space, and the
marked-root incidence are isomorphic presentations of one morphism.  The
identifications are not literal equalities of coordinate formulas.

The word ``exceptionally'' belongs to the assertion
$X_\tau\simeq\A^3$.  The target is visibly affine three-space, but the same
need not be expected of a simple-root incidence for an arbitrary slice.  The
type-$(2,1)$ tensor is the special geometric datum behind the construction.

\section{Resultant gauge and marked roots}

The volume form gives an intrinsic scalar normalization of the
linear--quadratic resultant.  For $\ell\in W^\vee$, let
$v_\ell\in W$ be determined by

\[
  \iota_{v_\ell}\omega=\ell.
\]

For $q\in\Sym^2W^\vee$, define

\[
  \Res_\omega(\ell,q)=q(v_\ell).
\]

In an $\omega$-unimodular basis this is the classical binary resultant of a
linear and a quadratic form.  In particular, it is nonzero precisely when
$\ell$ and $q$ are coprime.  Define the normalized factor space

\begin{equation}\label{eq:Y}
  Y_{\omega,\tau}=
  \left\{(\ell,q)\in W^\vee\times\Sym^2W^\vee:
  \Res_\omega(\ell,q)=1,\quad
  \langle\tau,\ell q\rangle=1\right\}.
\end{equation}

\begin{lemma}[Resultant gauge]\label{lem:gauge}
Multiplication and passage to the root line induce an isomorphism
\[
  Y_{\omega,\tau}\xrightarrow{\sim}X_\tau,
  \qquad
  (\ell,q)\longmapsto(\ell q,\PP(\ker\ell)).
\]
\end{lemma}

\begin{proof}
A marked root of $C$ is a factorization $C=\ell q$, unique up to
\[
  (\ell,q)\longmapsto(s\ell,s^{-1}q),\qquad s\in\Gm.
\]
The root is simple if and only if $\ell$ and $q$ are coprime.  Moreover,
\[
  \Res_\omega(s\ell,s^{-1}q)
  =s\,\Res_\omega(\ell,q).
\]
Thus every coprime scaling orbit has a unique resultant-one representative.
Because the normalizing scalar is the inverse of a regular nowhere-zero
function, this construction is algebraic in both directions.
\end{proof}

The weight one in the last formula is decisive: multiplication has a
$\Gm$-gauge, and the resultant cuts each coprime orbit exactly once.

There is also a root-gradient formulation.  If $C=\ell q$ and
$r=v_\ell$, then

\[
  \dd C_r=q(r)\ell
  =\Res_\omega(\ell,q)\,\iota_r\omega.
\]

Consequently $Y_{\omega,\tau}$ is equivalently the space of pairs

\begin{equation}\label{eq:gradient-normalization}
  (C,r)\in H_\tau\times(W\setminus\{0\}),
  \qquad \dd C_r=\iota_r\omega.
\end{equation}

Euler's identity gives $3C(r)=\dd C_r(r)=0$, while the right-hand side is a
nonzero covector.  Thus the equation simultaneously says that $r$ is a root
and that it is simple.  Conversely, every simple projective root has a unique
representative satisfying \eqref{eq:gradient-normalization}.

\section{A coordinate-free proof of etaleness}

Consider the five-dimensional polynomial map

\[
  \Psi\colon W^\vee\times\Sym^2W^\vee
  \longrightarrow\Sym^3W^\vee\times\A^1,
  \qquad
  (\ell,q)\longmapsto
  \bigl(\ell q,\Res_\omega(\ell,q)\bigr).
\]

\begin{proposition}\label{prop:etale}
The morphism $\Psi$ is etale wherever $\Res_\omega(\ell,q)\ne0$.
Consequently, $\pi\colon X_\tau\to H_\tau$ is etale.
\end{proposition}

\begin{proof}
Suppose that $\ell$ and $q$ are coprime.  If
\[
  \dot\ell q+\ell\dot q=0,
\]
then unique factorization implies that $\ell$ divides $\dot\ell$.
Degree considerations give
$(\dot\ell,\dot q)=a(\ell,-q)$ for some $a\in\C$.  Hence the kernel of the
differential of multiplication is exactly the infinitesimal gauge direction.

The resultant has gauge weight one, so
\[
  \dd\Res_\omega\big|_{(\ell,q)}(\ell,-q)
  =\Res_\omega(\ell,q)\ne0.
\]
Therefore $\dd\Psi$ is injective, hence an isomorphism between tangent spaces
of dimension five.  Finally,
\[
  Y_{\omega,\tau}
  =\Psi^{-1}(H_\tau\times\{1\}),
\]
and its map to $H_\tau$ is the base change of $\Psi$ along
$H_\tau\times\{1\}$.  Etaleness is preserved by base change, and
Lemma~\ref{lem:gauge} identifies this map with $\pi$.
\end{proof}

In coordinates the same argument is the relative-invariant identity

\[
  \det \dd\bigl(\operatorname{coeff}(\ell q),
  \Res_\omega(\ell,q)\bigr)
  =\pm\Res_\omega(\ell,q)^2.
\]

The differential proof explains the cancellation without expanding a
$5\times5$ determinant.

\section{The affine-space miracle}\label{sec:affine-chart}

It remains to prove the substantive assertion $X_\tau\simeq\A^3$.  Choose an
adapted basis $T,S$ of $W^\vee$, normalized so that
$\omega=T\wedge S$ and $\langle\tau,C\rangle$ extracts the coefficient of
$T^2S$.  Write

\[
  \ell=xT+\beta S,
  \qquad
  q=\gamma T^2+\delta TS+\epsilon S^2.
\]

Then \eqref{eq:Y} becomes the two equations

\begin{equation}\label{eq:slice}
  x\delta+\beta\gamma=1,
  \qquad
  x^2\epsilon-x\beta\delta+\beta^2\gamma=1.
\end{equation}

This resultant-normalized affine presentation and the chart below are the
factorization construction given by Lou \cite{Lou}; we include the inverse to
make the global affine-space assertion explicit.

\begin{proposition}[Global affine chart]\label{prop:chart}
The map $\chi\colon\A^3_{x,y,z}\to Y_{\omega,\tau}$ defined by
\begin{align*}
  \beta&=1+xy,\\
  \gamma&=\frac{2-3xy+x^2z}{2},\\
  \delta&=-\frac{xz-y+x^2yz-3xy^2}{2},\\
  \epsilon&=-z(1+xy)^2+4y^2+3xy^3
\end{align*}
is an isomorphism.  Its polynomial inverse is
\begin{equation}\label{eq:chart-inverse}
  x=x,\qquad
  y=2\beta\delta-x\epsilon,
  \qquad
  z=4y\delta+2\gamma y^2-\epsilon.
\end{equation}
\end{proposition}

\begin{proof}
Substitution verifies both equations in \eqref{eq:slice}.  Conversely, the
two equations of the slice imply
\[
  \beta-1=x(2\beta\delta-x\epsilon)=xy,
\]
so $\beta=1+xy$.  Substitution into the first slice equation, followed by
the definition of $z$ in \eqref{eq:chart-inverse}, gives
\[
  2\gamma=2-3xy+x^2z
\]
and then
\[
  2\delta=y(1+3xy)-xz(1+xy).
\]
Finally the formula defining $z$ gives
\[
  \epsilon=-z(1+xy)^2+y^2(4+3xy).
\]
These are precisely the displayed coordinate functions.  Both compositions
are therefore the identity, including on $x=0$; no divisor has been removed.
\end{proof}

This completes the proof that both source and target of $\pi$ are affine
three-spaces.  Together with Proposition~\ref{prop:etale}, it proves the
etaleness assertion of Theorem~\ref{thm:main}.

\section{The explicit Keller map}

Write a cubic in $H_\tau$ as

\[
  C=c_3T^3+T^2S+c_1TS^2+c_0S^3
\]

and use affine target coordinates $(2c_3,2c_1,c_0)$.  Multiplication
$C=\ell q$, followed by the chart of Proposition~\ref{prop:chart}, yields
$G=(G_1,G_2,G_3)$ with

\begin{align*}
G_1={}&x^3z-3x^2y+2x,\\
G_2={}&-3x^3y^2z+9x^2y^3-6x^2yz+12xy^2-3xz+y,\\
G_3={}&-x^3y^3z+3x^2y^4-3x^2y^2z+7xy^3-3xyz+4y^2-z.
\end{align*}

Because $G$ is etale, its Jacobian determinant is a unit of
$\C[x,y,z]$, hence a nonzero constant.  Its linear part at the origin is

\[
  (x,y,z)\longmapsto(2x,y,-z),
\]

and therefore

\[
  \boxed{\det JG=-2.}
\]

The map announced by Alp\"oge \cite{Alpoge}, in the original coordinate order,
is

\begin{align*}
F_1&=(1+xy)^3z+y^2(1+xy)(4+3xy),\\
F_2&=y+3x(1+xy)^2z+3xy^2(4+3xy),\\
F_3&=2x-3x^2y-x^3z.
\end{align*}

Indeed, if $S_0(x,y,z)=(x,y,-z)$ and
$R_0(a,b,c)=(c,b,a)$, then

\[
  G=R_0\circ F\circ S_0.
\]

Thus $F$ and $G$ are exactly left--right linearly equivalent presentations of
the same etale morphism.  This is an isomorphism of source--target arrows, not
strict dynamical conjugacy: the source and target coordinate changes are
independent.

There is an equally compact normalized-root certificate.  For
$F(x,y,z)=(A,B,C)$ define

\[
  Q_{A,B,C}(U,V)=CU^3-2U^2V+BUV^2-2AV^3
\]

and set $r=(U,V)=(1+xy,x)$.  Directly,

\[
  Q_{A,B,C}(r)=0,
  \qquad
  (Q_{A,B,C})_U(r)=2V,
  \qquad
  (Q_{A,B,C})_V(r)=-2U.
\]

For $\omega_0=\dd U\wedge\dd V$, this says

\[
  \dd Q_{A,B,C}\big|_r=-2\,\iota_r\omega_0.
\]

It is the same marked-root normalization as
\eqref{eq:gradient-normalization}, after the harmless source and target
rescalings relating $F$ and $G$.

\section{Fibers, discriminant, and infinity}

The fiber description is immediate from the moduli interpretation.

\begin{proposition}\label{prop:fibers}
For $C\in H_\tau$, the geometric fiber $\pi^{-1}(C)$ is the reduced set of
simple projective roots of $C$.  Its cardinality is
\[
  \#\pi^{-1}(C)=
  \begin{cases}
  3,&C\text{ has three simple roots},\\
  1,&C\text{ has a double root and a simple root},\\
  0,&C\text{ has a triple root}.
  \end{cases}
\]
In particular, $\pi$ has geometric degree three and is not injective.
\end{proposition}

For instance, $C_0=TS(T+S)$ lies in $H_\tau$ and has three possible marked
linear factors.  In the coordinates of $G$ these give

\[
  G(-1,2,-8)=G(0,2,16)=G(1,-1,-5)=(0,2,0).
\]

Equivalently, in the original coordinates,

\[
\begin{aligned}
F\left(0,0,-\frac14\right)
&=F\left(1,-\frac32,\frac{13}{2}\right)\\
&=F\left(-1,\frac32,\frac{13}{2}\right)
=\left(-\frac14,0,0\right).
\end{aligned}
\]

To see the boundary mechanism, introduce the full incidence

\[
  \overline X_\tau=
  \left\{(C,[r])\in H_\tau\times\PP(W):C(r)=0\right\}.
\]

Its projection to $H_\tau$ is finite flat of degree three.  Its ramification
locus consists of repeated marked roots, and its branch locus is the binary
cubic discriminant.  The source $X_\tau$ is obtained by deleting precisely
that ramification locus upstairs.  Over the complement of the discriminant,
no point is deleted and $\pi$ is finite etale.  Over a double-root cubic the
length-two ramified point is removed while the simple-root point remains; at
a triple-root cubic every point is removed.

In the affine chart $X_\tau\simeq\A^3$, the deleted points appear at infinity.
Thus $\pi$ is etale and quasi-finite but not finite or proper.  It should not
be called a topological three-sheeted covering.  Its generic monodromy is the
ordinary $S_3$-monodromy of the roots of a cubic.

For the original target coordinates $(a,b,c)$, a normalized discriminant
(one quarter of the standard binary-cubic discriminant) is

\[
  \Delta(a,b,c)
  =b^2-cb^3-16a+18abc-27a^2c^2.
\]

The nonproperness set is exactly $V(\Delta)$.  Triple-root cubics form the
omitted curve

\[
  (a,b,c)=
  \left(\frac{\rho^2}{3},2\rho,\frac{2}{3\rho}\right),
  \qquad \rho\in\C^*.
\]

This is the geometric reason the map can be everywhere locally invertible
without being globally invertible: its missing inverse branches do not
ramify at finite source points; they escape through the boundary of the
affine chart.

\section{Conclusion}

The coordinate formula conceals a short moduli problem:

\[
  \boxed{
  \text{binary cubic in a tangential affine slice}
  +\text{ one normalized simple root}
  \ \longmapsto\ \text{the cubic}.}
\]

The resultant-one condition removes the factor-rescaling gauge, and its
weight one proves etaleness.  A general cubic has three roots, which explains
the degree.  The exceptional global chart
$X_\tau\simeq\A^3$ converts this forgetful morphism into a polynomial
self-map of affine three-space.  The discriminant records exactly where
inverse branches leave that chart.  In this form, constant Jacobian,
noninjectivity, and nonproperness are three aspects of the same marked-root
geometry.

\begin{thebibliography}{9}

\bibitem{Alpoge}
L.~Alp\"oge,
\newblock X post announcing the explicit three-dimensional polynomial map,
\newblock 20 July 2026, 02:19:17 UTC,
\newblock \url{https://x.com/__alpoge__/status/2079028340955197566}.

\bibitem{Jiang}
A.~Jiang (@davikrehalt),
\newblock X post giving the projective symmetric-product marked-root
  formulation,
\newblock 20 July 2026, 12:02:19 UTC,
\newblock \url{https://x.com/davikrehalt/status/2079175065695035442}.

\bibitem{Lou}
A.~Lou,
\newblock Deriving an explicit polynomial counterexample to the Jacobian
  conjecture: a reproducible cubic-factor and resultant construction,
\newblock research note, 20 July 2026,
\newblock \url{https://aaronlou.com/jacobian_counterexample_derivation.pdf}.

\bibitem{BCW}
H.~Bass, E.~H. Connell, and D.~Wright,
\newblock The Jacobian conjecture: reduction of degree and formal expansion
  of the inverse,
\newblock \emph{Bull. Amer. Math. Soc. (N.S.)} \textbf{7} (1982), 287--330.

\bibitem{Essen}
A.~van den Essen,
\newblock \emph{Polynomial Automorphisms and the Jacobian Conjecture},
\newblock Progress in Mathematics 190, Birkh\"auser, 2000.

\end{thebibliography}

\end{document}
